% !TEX TS-program = pdflatex
% Compilation : pdflatex rapport.tex (×2) depuis le dossier pfe/
\documentclass[12pt,a4paper,openany]{report}

\usepackage[utf8]{inputenc}
\usepackage[T1]{fontenc}
\usepackage{lmodern}
\usepackage[french]{babel}
\usepackage{csquotes}
\usepackage{microtype}

\usepackage[a4paper,top=2.5cm,bottom=2.5cm,left=3cm,right=2.5cm]{geometry}
\usepackage{setspace}
\onehalfspacing

\usepackage{xcolor}
\definecolor{mainblue}{RGB}{41,128,185}
\definecolor{maindark}{RGB}{44,62,80}

\usepackage{titlesec}
\titleformat{\chapter}[display]
  {\normalfont\huge\bfseries\color{mainblue}}
  {\chaptertitlename\ \thechapter}{20pt}{\Huge}

\usepackage{fancyhdr}
\setlength{\headheight}{14.5pt}
\pagestyle{fancy}
\fancyhf{}
\fancyhead[L]{\small\leftmark}
\fancyhead[R]{\small MoroccoCheck}
\fancyfoot[C]{\thepage}
\renewcommand{\headrulewidth}{0.4pt}
\fancypagestyle{plain}{
  \fancyhf{}
  \fancyfoot[C]{\thepage}
  \renewcommand{\headrulewidth}{0pt}
}

\usepackage{float}
\usepackage{graphicx}
\graphicspath{{figures/}}

\usepackage{hyperref}
\hypersetup{
  colorlinks=true,
  linkcolor=mainblue,
  citecolor=mainblue,
  urlcolor=maindark,
  pdftitle={MoroccoCheck - Rapport PFE},
  pdfauthor={BADDOU, EL IDRISSI}
}

\usepackage{listings}
\lstset{
  basicstyle=\ttfamily\small,
  breaklines=true,
  frame=single,
  rulecolor=\color{maindark},
  numbers=left,
  numberstyle=\tiny\color{gray},
  keywordstyle=\color{mainblue},
  commentstyle=\itshape\color{gray},
  captionpos=b
}

\usepackage{booktabs}
\usepackage{array}
\usepackage{longtable}

\newcommand{\placeholderfig}[2]{%
  \fbox{\parbox[c][0.18\textheight]{\linewidth}{\centering\vfill\textit{#1}\\\small\texttt{#2}\vfill}}%
}

\usepackage{enumitem}
\setlist{itemsep=2pt,parsep=0pt}

\begin{document}

\begin{titlepage}
\centering
\vspace*{1cm}

% Logo (optionnel)
\IfFileExists{figures/logo_est.png}{%
  \includegraphics[width=4cm]{figures/logo_est.png}\\[1.5cm]
}{%
  \fbox{\parbox[c][2.5cm]{8cm}{\centering\textit{[Logo établissement — placer \texttt{figures/logo\_est.png}]}}}\\[1.5cm]
}

{\Large École supérieure de technologie -- Agadir\par}
\vspace{0.8cm}
{\large Filière : Conception et développement des logiciels\par}
\vspace{2cm}

\rule{\textwidth}{1.5pt}\\[0.6cm]
{\LARGE\bfseries\color{mainblue} MoroccoCheck\par}
\vspace{0.4cm}
{\large Plateforme touristique communautaire vérifiée par géolocalisation\par}
\rule{\textwidth}{1.5pt}

\vfill

{\large\textbf{Rapport de Projet de Fin d'Études (PFE)}\par}
\vspace{1.5cm}

\begin{tabular}{rl}
\textbf{Réalisé par :} & Abderrahmane \textsc{Baddou}\\
                      & Ahmed \textsc{El Idrissi}\\[1em]
\textbf{Encadré par :} & Pr. Aïcha \textsc{Bakki}\\[1em]
\textbf{Devant le jury :} & Pr. Abderrahim \textsc{Sabour}\\
                         & Pr. Aïcha \textsc{Bakki}\\
\end{tabular}

\vfill

{\large Année universitaire 2025--2026\par}

\end{titlepage}

\cleardoublepage
\pagenumbering{roman}
\chapter*{Dédicace}
\addcontentsline{toc}{chapter}{Dédicace}

\noindent\textit{[Texte de dédicace à personnaliser.]}

\vfill

\chapter*{Remerciements}
\addcontentsline{toc}{chapter}{Remerciements}

Nous tenons à exprimer notre gratitude à \textbf{Pr. Aïcha Bakki}, notre encadrante pédagogique, pour son accompagnement, ses conseils et la qualité de son suivi tout au long de la réalisation de ce projet.

Nous remercions également les membres du jury, \textbf{Pr. Abderrahim Sabour} et \textbf{Pr. Aïcha Bakki}, pour le temps qu'ils consacreront à l'évaluation de ce travail.

Enfin, nous adressons nos remerciements à l'ensemble du corps enseignant de l'\textbf{École supérieure de technologie d'Agadir} pour la formation dispensée.

\vfill

\chapter*{Résumé}
\addcontentsline{toc}{chapter}{Résumé}

\noindent
\textbf{MoroccoCheck} est une plateforme touristique \textbf{communautaire} dédiée au Maroc, composée d'une \textbf{API REST} (\textbf{Node.js} / \textbf{Express} / \textbf{MySQL}), d'une application mobile multiplateforme (\textbf{Flutter}) et d'une interface d'administration web (\textbf{React} / \textbf{Vite}). Le système permet de consulter des sites touristiques, d'effectuer des \textbf{check-ins validés par géolocalisation} (formule de Haversine, contrôle de rayon et de précision), de publier des avis et de mettre en œuvre une \textbf{gamification} (points, badges, classement). Des rôles métiers (\textsc{Tourist}, \textsc{Contributor}, \textsc{Professional}, \textsc{Admin}) encadrent les usages ; les administrateurs assurent la \textbf{modération} des contenus via l'interface web.

\vspace{1em}
\noindent\textbf{Mots-clés :} tourisme numérique ; géolocalisation ; API REST ; Flutter ; gamification ; modération ; JWT ; MySQL.

\vspace{2em}
\chapter*{Abstract}
\addcontentsline{toc}{chapter}{Abstract}

\noindent
\textbf{MoroccoCheck} is a community \textbf{tourism platform} for Morocco, consisting of a \textbf{REST API} (\textbf{Node.js} / \textbf{Express} / \textbf{MySQL}), a cross-platform mobile app (\textbf{Flutter}), and a web admin interface (\textbf{React} / \textbf{Vite}). The system supports discovery of tourist sites, \textbf{GPS-validated check-ins} (Haversine distance, radius and accuracy checks), reviews, and \textbf{gamification} (points, badges, leaderboard). Business roles (\textsc{Tourist}, \textsc{Contributor}, \textsc{Professional}, \textsc{Admin}) structure access rights; administrators moderate content through the web UI.

\vspace{1em}
\noindent\textbf{Keywords:} digital tourism; geolocation; REST API; Flutter; gamification; moderation; JWT; MySQL.

\tableofcontents
\clearpage
\listoffigures
\clearpage
\listoftables
\clearpage
\chapter*{Liste des abréviations}
\addcontentsline{toc}{chapter}{Liste des abréviations}

\begin{tabular}{@{}p{3.2cm}p{11cm}@{}}
API & Application Programming Interface \\
CRUD & Create, Read, Update, Delete \\
CI/CD & Intégration continue / déploiement continu \\
GPS & Global Positioning System \\
HTTP & HyperText Transfer Protocol \\
HTTPS & HTTP sécurisé (TLS) \\
JSON & JavaScript Object Notation \\
JWT & JSON Web Token \\
MySQL & Système de gestion de base de données relationnelle \\
OAuth & Protocole d'autorisation ouvert \\
ORM & Object-Relational Mapping \\
PFE & Projet de fin d'études \\
REST & Representational State Transfer \\
SPA & Single Page Application \\
SQL & Structured Query Language \\
UML & Unified Modeling Language \\
\end{tabular}

\clearpage


\cleardoublepage
\pagenumbering{arabic}
\setcounter{page}{1}

\chapter*{Introduction générale}
\addcontentsline{toc}{chapter}{Introduction générale}
\markboth{Introduction générale}{Introduction générale}

Le secteur du \textbf{tourisme} s'appuie de plus en plus sur le \textbf{numérique} pour informer, orienter et fidéliser les visiteurs. Les voyageurs consultent des avis, des cartes et des applications mobiles ; pourtant, l'information disponible est souvent \textbf{hétérogène}, \textbf{peu actualisée} ou \textbf{difficile à vérifier} sur le terrain.

Le projet \textbf{MoroccoCheck} répond à ce constat en proposant une \textbf{plateforme touristique communautaire} centrée sur le Maroc. L'application combine la consultation de \textbf{sites touristiques}, la \textbf{vérification de présence} par \textbf{géolocalisation} (check-in GPS avec contrôle de distance), la publication d'\textbf{avis} et une \textbf{gamification} (points, badges, classement). Un rôle \textbf{professionnel} permet aux gestionnaires de lieux de revendiquer des fiches et d'interagir avec les retours visiteurs, tandis que les \textbf{administrateurs} assurent la \textbf{modération} et le pilotage via une interface web dédiée.

\section{Problématique}

La question centrale de ce mémoire est la suivante :

\begin{quote}
\textit{Comment concevoir une plateforme touristique communautaire capable d'améliorer la fiabilité des informations grâce à la géolocalisation, à la participation des utilisateurs et à la modération, tout en restant évolutive sur le plan technique ?}
\end{quote}

\section{Objectifs du projet}

Les objectifs principaux sont :
\begin{itemize}[leftmargin=*]
  \item fournir un \textbf{catalogue consultable} de sites (liste, détail, catégories, carte) ;
  \item mettre en œuvre des \textbf{check-ins validés par GPS} (distance, précision, règles métier) ;
  \item permettre la \textbf{contribution communautaire} (avis, notes, médias selon les parcours) ;
  \item encourager l'engagement par la \textbf{gamification} ;
  \item gérer les \textbf{rôles} (touriste, contributeur, professionnel, administrateur) et les demandes d'évolution de statut ;
  \item offrir un \textbf{back-office} pour statistiques et modération.
\end{itemize}

\section{Méthodologie et structure du document}

Le travail s'appuie sur une analyse des besoins, une phase de conception (modélisation UML, base de données), puis une réalisation incrémentale (API, mobile, administration) avec tests et intégration continue.

Ce rapport est organisé en \textbf{trois chapitres} : le \textbf{cadre général et l'analyse des besoins}, la \textbf{conception de la solution}, puis la \textbf{réalisation et l'implémentation}. Une \textbf{conclusion générale}, un chapitre sur les \textbf{outils de développement} et une \textbf{bibliographie} complètent le document.

\chapter{Cadre général et analyse des besoins}

\section{Contexte et motivation}

Les touristes et résidents qui souhaitent découvrir le Maroc se heurtent fréquemment à des informations \textbf{fragmentaires} ou \textbf{peu fiables}. Les guides généralistes et les réseaux sociaux offrent une visibilité utile mais rarement une \textbf{validation terrain} structurée. \textbf{MoroccoCheck} vise à \textbf{centraliser} la découverte des sites, à \textbf{ancrer les retours utilisateurs dans une preuve de présence} (géolocalisation) et à \textbf{encadrer la qualité} par la modération.

\section{Acteurs et périmètre fonctionnel}

Le tableau~\ref{tab:acteurs} résume les acteurs et leurs principales interactions avec le système.

\begin{table}[H]
\centering
\caption{Acteurs métier et périmètre fonctionnel}
\label{tab:acteurs}
\begin{tabular}{@{}p{3.2cm}p{10.5cm}@{}}
\toprule
\textbf{Acteur} & \textbf{Fonctionnalités principales} \\
\midrule
Touriste & Inscription, connexion (email / Google), consultation des sites, carte, check-in GPS, avis, profil, demande de passage contributeur. \\
Contributeur & Périmètre élargi de contribution communautaire (hérite des usages touriste selon le modèle de rôles). \\
Professionnel & Création ou mise à jour de fiches, revendication de site, espace professionnel, réponse aux avis. \\
Administrateur & Tableaux de bord, modération des sites et des avis, gestion des utilisateurs, traitement des demandes contributeur. \\
\bottomrule
\end{tabular}
\end{table}

\section{Besoins fonctionnels}

Les besoins suivants structurent l'application :

\begin{itemize}[leftmargin=*]
  \item \textbf{Authentification et comptes :} inscription, connexion, rafraîchissement de jetons, profil, sessions persistées.
  \item \textbf{Catalogue :} catégories, liste et détail de sites, médias, carte.
  \item \textbf{Check-in GPS :} envoi des coordonnées, calcul de distance (Haversine), contrôle de rayon et de précision, enregistrement du check-in.
  \item \textbf{Avis :} création, modification, modération, réponse éventuelle du professionnel.
  \item \textbf{Gamification :} points, badges, statistiques, classement.
  \item \textbf{Administration :} statistiques globales, files de modération, gestion des comptes.
\end{itemize}

\section{Besoins non fonctionnels}

\begin{itemize}[leftmargin=*]
  \item \textbf{Sécurité :} authentification par JWT, mots de passe hashés, validation des entrées, contrôle d'accès par rôle.
  \item \textbf{Performance :} indexation SQL, vues pour agrégats, limitation de débit optionnelle (Redis ou mémoire).
  \item \textbf{Disponibilité :} points de contrôle \texttt{/api/health}, journalisation des requêtes, suivi d'erreurs (Sentry).
  \item \textbf{Évolutivité :} architecture en couches, migrations SQL, API REST documentée par le code et les README.
\end{itemize}

\section{État de l'art succinct}

Le tableau~\ref{tab:concurrents} positionne le projet par rapport à des solutions connues (à affiner selon vos sources bibliographiques).

\begin{table}[H]
\centering
\caption{Comparaison qualitative avec des solutions existantes}
\label{tab:concurrents}
\small
\begin{tabular}{@{}p{3cm}p{5cm}p{5cm}@{}}
\toprule
\textbf{Solution} & \textbf{Forces} & \textbf{Limites} \\
\midrule
TripAdvisor & Large base d'avis, notoriété & Peu de preuve de présence sur site \\
Google Maps & Cartographie, avis & Générique, peu spécifique au positionnement métier MoroccoCheck \\
Guides / blogs & Contenu éditorial & Actualisation irrégulière, pas de validation GPS intégrée \\
\bottomrule
\end{tabular}
\end{table}

MoroccoCheck se distingue par l'association explicite \textbf{avis / présence géolocalisée} et par un \textbf{dispositif de modération} et de \textbf{rôles métiers} intégré au même produit (mobile + API + admin).

\chapter{Conception de la solution}

\section{Architecture logique}

L'architecture adopte un modèle \textbf{client--serveur} à trois composants principaux :
\begin{itemize}[leftmargin=*]
  \item une \textbf{API REST} monolithique (\textbf{Node.js} / \textbf{Express}) ;
  \item une application \textbf{Flutter} (utilisateurs) ;
  \item une interface \textbf{React} / \textbf{Vite} (administration).
\end{itemize}

La logique serveur suit une organisation \textbf{Routes $\rightarrow$ Contrôleurs $\rightarrow$ Services}, avec des \textbf{middlewares} (authentification JWT, gestion d'erreurs, limitation de débit, téléversement). Les données sont stockées dans \textbf{MySQL} (moteur InnoDB, encodage \texttt{utf8mb4}).

\section{Diagramme de composants}

La figure~\ref{fig:components} présente la répartition en couches et les dépendances externes (Firebase, OAuth Google, fournisseur de cartographie, CI GitHub Actions).

\begin{figure}[H]
\centering
\IfFileExists{figures/component_architecture_moroccocheck.png}{%
  \includegraphics[width=\linewidth]{figures/component_architecture_moroccocheck.png}%
}{%
  \placeholderfig{Diagramme de composants — architecture}{figures/component\_architecture\_moroccocheck.png}%
}
\caption{Diagramme de composants — architecture MoroccoCheck}
\label{fig:components}
\end{figure}

\section{Modèle de données}

Les entités centrales incluent notamment : \textbf{User}, \textbf{TouristSite}, \textbf{Category}, \textbf{CheckIn}, \textbf{Review}, \textbf{Badge}, \textbf{UserBadge}, \textbf{ContributorRequest}, ainsi que des tables de liaison et de session. Les \textbf{vues SQL} (ex.\ \texttt{v\_user\_stats}, \texttt{v\_site\_details}) facilitent les statistiques ; des \textbf{déclencheurs} maintiennent la cohérence des compteurs et des moyennes de notes lors des insertions ou suppressions d'avis et de check-ins.

\section{Diagramme de classes}

Le diagramme de classes (figure~\ref{fig:classes}) synthétise les associations principales entre entités du domaine.

\begin{figure}[H]
\centering
\IfFileExists{figures/class_diagram_moroccocheck.png}{%
  \includegraphics[width=\linewidth]{figures/class_diagram_moroccocheck.png}%
}{%
  \placeholderfig{Diagramme de classes du domaine}{figures/class\_diagram\_moroccocheck.png}%
}
\caption{Diagramme de classes — domaine métier MoroccoCheck}
\label{fig:classes}
\end{figure}

\section{Cas d'utilisation}

La figure~\ref{fig:uc} illustre les cas d'utilisation globaux et la généralisation entre touriste et contributeur.

\begin{figure}[H]
\centering
\IfFileExists{figures/use_case_global_moroccocheck.png}{%
  \includegraphics[width=\linewidth]{figures/use_case_global_moroccocheck.png}%
}{%
  \placeholderfig{Diagramme de cas d'utilisation global}{figures/use\_case\_global\_moroccocheck.png}%
}
\caption{Diagramme de cas d'utilisation global}
\label{fig:uc}
\end{figure}

\section{Scénarios dynamiques}

\subsection{Authentification}

La figure~\ref{fig:seq-auth} modélise le scénario de connexion : validation des entrées, recherche utilisateur, comparaison de mot de passe (\texttt{bcrypt}), émission de jetons JWT et persistance de session.

\begin{figure}[H]
\centering
\IfFileExists{figures/sequence_auth_moroccocheck.png}{%
  \includegraphics[width=0.95\linewidth]{figures/sequence_auth_moroccocheck.png}%
}{%
  \placeholderfig{Séquence authentification}{figures/sequence\_auth\_moroccocheck.png}%
}
\caption{Diagramme de séquence — authentification (login)}
\label{fig:seq-auth}
\end{figure}

\subsection{Check-in GPS}

La figure~\ref{fig:seq-checkin} décrit le scénario de validation d'un check-in : acquisition GPS côté mobile, appel API authentifié, calcul de distance, règles métier, persistance et mise à jour de la gamification.

\begin{figure}[H]
\centering
\IfFileExists{figures/sequence_checkin_gps_moroccocheck.png}{%
  \includegraphics[width=0.95\linewidth]{figures/sequence_checkin_gps_moroccocheck.png}%
}{%
  \placeholderfig{Séquence check-in GPS}{figures/sequence\_checkin\_gps\_moroccocheck.png}%
}
\caption{Diagramme de séquence — check-in GPS}
\label{fig:seq-checkin}
\end{figure}

\section{Activité — gamification}

Le traitement des points, du rang et des badges est schématisé figure~\ref{fig:act-gamif}. Les seuils de rang dans l'implémentation sont paramétrés dans le code (\texttt{RANK\_THRESHOLDS}) ; le rapport académique peut les aligner sur cette source lors de la rédaction finale.

\begin{figure}[H]
\centering
\IfFileExists{figures/activity_gamification_moroccocheck.png}{%
  \includegraphics[width=\linewidth]{figures/activity_gamification_moroccocheck.png}%
}{%
  \placeholderfig{Diagramme d'activité — gamification}{figures/activity\_gamification\_moroccocheck.png}%
}
\caption{Diagramme d'activité — gamification (points, rangs, badges)}
\label{fig:act-gamif}
\end{figure}

\chapter{Réalisation et implémentation}

\section{Structure du dépôt logiciel}

Le dépôt regroupe trois applications principales :

\begin{verbatim}
App_Touriste/
  back-end/     # API Express (Node.js), scripts SQL, tests
  front-end/    # Application Flutter
  admin-web/    # Interface React / Vite
  .github/workflows/  # Intégration continue
\end{verbatim}

\section{Couche serveur (API)}

Le fichier \texttt{server.js} instancie l'application Express : sécurisation des en-têtes (Helmet), CORS, corps JSON, journalisation (Morgan), exposition des fichiers téléversés sous \texttt{/uploads}, enregistrement des routes \texttt{/api/...} et middleware d'erreur en fin de chaîne.

Les modules métier incluent notamment l'authentification (\texttt{/api/auth}), les sites (\texttt{/api/sites}), les check-ins (\texttt{/api/checkins}), les avis (\texttt{/api/reviews}), les utilisateurs et la gamification, ainsi que l'administration (\texttt{/api/admin}).

\subsection{Exemple : distance GPS (Haversine)}

L'extrait suivant illustre le calcul de distance utilisé pour valider un check-in (fichier \texttt{gps.utils.js}) :

\begin{figure}[H]
\begin{lstlisting}[language=Java]
export function calculateDistance(lat1, lon1, lat2, lon2) {
  const R = 6371000;
  const lat1Rad = (lat1 * Math.PI) / 180;
  const lon1Rad = (lon1 * Math.PI) / 180;
  const lat2Rad = (lat2 * Math.PI) / 180;
  const lon2Rad = (lon2 * Math.PI) / 180;
  const dLat = lat2Rad - lat1Rad;
  const dLon = lon2Rad - lon1Rad;
  const a =
    Math.sin(dLat / 2) ** 2 +
    Math.cos(lat1Rad) * Math.cos(lat2Rad) * Math.sin(dLon / 2) ** 2;
  const c = 2 * Math.atan2(Math.sqrt(a), Math.sqrt(1 - a));
  return R * c;
}
\end{lstlisting}
\caption{Calcul de distance Haversine (\texttt{back-end/src/utils/gps.utils.js})}
\label{fig:code-haversine}
\end{figure}

La fonction \texttt{isWithinRadius} compare ensuite cette distance aux seuils définis dans \texttt{GPS\_VALIDATION} (\texttt{constants.js}).

\section{Application mobile (Flutter)}

Le client Flutter s'organise par \textbf{fonctionnalités} (\texttt{lib/features/}) : authentification, sites, carte, profil, espace professionnel, paramètres. Le cœur technique (\texttt{lib/core/}) regroupe les constantes, le thème, le client HTTP (\texttt{api\_service.dart}), la géolocalisation, les notifications, la synchronisation hors ligne et le routeur (\texttt{go\_router}). L'état est géré notamment avec \textbf{Provider}.

\section{Interface d'administration}

L'administration est une \textbf{SPA} React servie par Vite, réservée aux comptes \textbf{ADMIN}. Elle consomme la même API via des requêtes HTTP JSON (URL de base configurable, ex.\ \texttt{VITE\_API\_BASE\_URL}).

\section{Sécurité}

Les principales mesures déployées incluent :
\begin{itemize}[leftmargin=*]
  \item authentification \textbf{JWT} avec contrôle des sessions en base ;
  \item hachage des mots de passe (\texttt{bcrypt}) ;
  \item validation des entrées (Joi, express-validator) ;
  \item limitation du débit configurable (mémoire ou Redis) ;
  \item séparation des droits selon le rôle utilisateur.
\end{itemize}

\section{Tests et qualité}

Le backend utilise \textbf{Mocha}, \textbf{Chai} et \textbf{Supertest}. Le dépôt définit un workflow \textbf{GitHub Actions} : installation des dépendances, base MySQL de test, migrations, exécution des tests ; build de l'admin ; analyse et compilation Flutter (APK de staging). L'analyse statique Flutter (\texttt{flutter analyze}) contribue à la qualité du code mobile.

\section{Limites et perspectives techniques}

Parmi les limites actuelles : déploiement production à finaliser, couverture de tests frontend à étendre, documentation API de type OpenAPI à industrialiser. Les perspectives incluent renforcement du monitoring, sauvegardes automatisées et évolutions fonctionnelles décrites dans la documentation projet (recommandations, tableaux de bord admin enrichis, etc.).

\section{Interfaces utilisateur}

Les captures d'écran des écrans principaux (connexion, accueil, liste et détail de site, carte, check-in, profil, espace pro, administration) sont à insérer dans une annexe ou dans ce chapitre une fois les figures exportées (PNG/PDF) sous \texttt{pfe/figures/}.

\chapter{Conclusion générale}

Ce projet de fin d'études a permis de concevoir et de réaliser \textbf{MoroccoCheck}, une plateforme touristique communautaire articulée autour d'une \textbf{API REST}, d'une application \textbf{Flutter} et d'une interface d'\textbf{administration web}. Les objectifs initiaux — catalogue de sites, \textbf{vérification par géolocalisation}, \textbf{avis}, \textbf{gamification}, rôles métiers et \textbf{modération} — sont couverts par une base logicielle cohérente et documentée.

Sur le plan technique, l'équipe a consolidé des compétences en développement \textbf{backend} (Node.js, MySQL, sécurité), \textbf{mobile} (Flutter, intégration API) et \textbf{frontend web} (React), ainsi qu'en modélisation (\textbf{UML}), gestion de version (\textbf{Git}) et intégration continue (\textbf{GitHub Actions}).

Les perspectives à moyen terme portent sur la généralisation du déploiement en environnement de production, l'élargissement des tests automatisés et l'enrichissement progressif des fonctionnalités (recommandations, analytics, expérience multilingue, etc.), dans la continuité des axes identifiés dans l'analyse du projet.

Ce travail démontre qu'une approche \textbf{itérative}, fondée sur une architecture claire et des besoins bien identifiés, permet de livrer un système complexe tout en gardant une trajectoire d'amélioration maîtrisée.

\chapter{Outils de développement}

\section{Environnement intégré et langages}

\begin{itemize}[leftmargin=*]
  \item \textbf{Édition :} Visual Studio Code, Android Studio ou environnement équivalent pour Flutter.
  \item \textbf{Langages :} JavaScript (Node.js), Dart (Flutter), TypeScript (admin web), SQL (MySQL).
\end{itemize}

\section{Gestion de version et collaboration}

\begin{itemize}[leftmargin=*]
  \item \textbf{Git} pour le suivi des versions ;
  \item hébergement sur \textbf{GitHub} (branche principale, pull requests, workflow CI).
\end{itemize}

\section{Outils de build et d'exécution}

\begin{itemize}[leftmargin=*]
  \item \textbf{npm} pour le backend et l'admin ;
  \item \textbf{Flutter SDK} (canal stable) pour le mobile ;
  \item \textbf{Vite} pour le bundling de l'interface administrateur.
\end{itemize}

\section{Qualité et supervision}

\begin{itemize}[leftmargin=*]
  \item analyse statique Flutter (\texttt{flutter analyze}) ;
  \item tests backend (Mocha) et pipeline CI ;
  \item \textbf{Sentry} pour le suivi d'erreurs côté API et mobile (configuration par environnement).
\end{itemize}

\section{Rédaction du rapport}

\begin{itemize}[leftmargin=*]
  \item \textbf{\LaTeX} (moteur \texttt{pdflatex}) pour la composition du présent document ;
  \item \textbf{PlantUML} pour les diagrammes UML exportés en PNG dans \texttt{pfe/figures/}.
\end{itemize}

\chapter{Bibliographie et webographie}

\begin{thebibliography}{99}

\bibitem{express}
OpenJS Foundation.
\newblock \textit{Express — web framework for Node.js}.
\newblock \url{https://expressjs.com/}, consulté en 2026.

\bibitem{flutter}
Google LLC.
\newblock \textit{Flutter documentation}.
\newblock \url{https://docs.flutter.dev/}, consulté en 2026.

\bibitem{mysql}
Oracle Corporation.
\newblock \textit{MySQL 8.0 Reference Manual}.
\newblock \url{https://dev.mysql.com/doc/refman/8.0/en/}, consulté en 2026.

\bibitem{jwt}
IETF.
\newblock RFC 7519 — JSON Web Token (JWT).
\newblock \url{https://www.rfc-editor.org/rfc/rfc7519}, 2015.

\bibitem{oauth}
IETF.
\newblock RFC 6749 — The OAuth 2.0 Authorization Framework.
\newblock \url{https://www.rfc-editor.org/rfc/rfc6749}, 2012.

\bibitem{react}
Meta Platforms, Inc.
\newblock \textit{React — A JavaScript library for building user interfaces}.
\newblock \url{https://react.dev/}, consulté en 2026.

\bibitem{vite}
Vite Team.
\newblock \textit{Vite — Next Generation Frontend Tooling}.
\newblock \url{https://vitejs.dev/}, consulté en 2026.

\bibitem{uml}
Object Management Group (OMG).
\newblock \textit{Unified Modeling Language}.
\newblock \url{https://www.omg.org/spec/UML/}, consulté en 2026.

\end{thebibliography}

\section*{Webographie complémentaire}
\addcontentsline{toc}{section}{Webographie complémentaire}

\begin{itemize}[leftmargin=*]
  \item Documentation Node.js : \url{https://nodejs.org/docs/}
  \item Documentation Dart : \url{https://dart.dev/guides}
  \item Référence PlantUML : \url{https://plantuml.com/}
  \item GitHub Actions : \url{https://docs.github.com/actions}
\end{itemize}


\end{document}
